\part{Part 1: Introduction to electrical machines and drives}
\title[Introduction]{Introduction to electrical machines and drives}  
\date{}  
\frame{\titlepage} 

%%%%%%%%%%%%%%%%%%%%%%%%%%%%%%%%%%%%%%%%%%%%%%%%%%%%%%%%%%%%%
%% What is an electrical machine ? %%
%%%%%%%%%%%%%%%%%%%%%%%%%%%%%%%%%%%%%%%%%%%%%%%%%%%%%%%%%%%%%
\begin{frame}
	\frametitle{What is an electrical machine ?}
	\begin{columns}
		\begin{column}{0.5\textwidth}
			\begin{defi}{Electrical machine}{elec_machine}
			   An electrical machine is a device that \hl{converts electrical energy into  mechanical energy} or vice versa.
			\end{defi}
			\vspace{0.25cm}
			\begin{itemize}
				\item Electrical energy is routed via machine's external wiring connected to the terminal box.
				\item Mechanical energy is transferred via the shaft (if it is a rotatory machine).
			\end{itemize}
		\end{column}
		\begin{column}{0.5\textwidth}
			\begin{figure}
				\centering
				\includegraphics[width=0.95\textwidth]{fig/lec01/Induction_machine_opened.pdf}
				\caption{Example of an electrical machine (source: derived from \href{https://commons.wikimedia.org/wiki/File:TMW_50906_Schnittmodell_einer_Drehstrommaschine_(Asynchronmaschine).jpg}{Wikimedia Commons}, public domain)}
			\end{figure}
		\end{column}
		\end{columns}
\end{frame}


%%%%%%%%%%%%%%%%%%%%%%%%%%%%%%%%%%%%%%%%%%%%%%%%%%%%%%%%%%%%%
%% What is an electrical drive ? %%
%%%%%%%%%%%%%%%%%%%%%%%%%%%%%%%%%%%%%%%%%%%%%%%%%%%%%%%%%%%%%
\begin{frame}
	\frametitle{What is an electrical drive ?}
	\begin{columns}
		\begin{column}{0.4\textwidth}
			\begin{defi}{Electrical drive}{elec_drive}
			   An electrical drive is a system that \hl{controls the torque, speed or position of an electrical machine} connected to some mechanical process.
			\end{defi}
			\vspace{0.25cm}
			\begin{itemize}
				\item Integrates the 'stupid' electrical machine into an 'intelligent' controlled system.
				\item The energy source and mechanical process ('load') are not part of the drive system.
			\end{itemize}
		\end{column}
		\begin{column}{0.6\textwidth}
			\begin{figure}
				\centering
				\includegraphics[width=0.95\textwidth]{fig/lec01/Electrical_Drive_Block_Overview.pdf}
				\caption{Block diagram of an electrical drive}
			\end{figure}
		\end{column}
	\end{columns}
\end{frame}

%%%%%%%%%%%%%%%%%%%%%%%%%%%%%%%%%%%%%%%%%%%%%%%%%%%%%%%%%%%%%
%% Examples of electrical machine and drive applications (1) %%
%%%%%%%%%%%%%%%%%%%%%%%%%%%%%%%%%%%%%%%%%%%%%%%%%%%%%%%%%%%%%
\begin{frame}
	\frametitle{Examples of electrical machine and drive applications (1)}
	\begin{figure}
		\centering
		\begin{subfigure}[b]{0.49\textwidth}
			\centering
			\includegraphics[width=0.5\textwidth]{fig/lec01/Electric_Car_recharging.jpg}
			\caption{Electric cars (source: \href{https://commons.wikimedia.org/wiki/File:Electric_Car_recharging.jpg}{Wikimedia Commons}, M.~Movchin and F. Mueller, \href{https://creativecommons.org/licenses/by-sa/3.0/deed.en}{CC BY-SA 3.0})}
		\end{subfigure}
		\hfill
		\begin{subfigure}[b]{0.49\textwidth}
			\centering
			\includegraphics[width=0.5\textwidth]{fig/lec01/sky-farm-windmill.jpg}
			\caption{Wind turbine generators (source: \href{https://pxhere.com/en/photo/954757}{pxhere.com}, public domain)}
		\end{subfigure}
		\\
		\begin{subfigure}[b]{0.49\textwidth}
			\centering
			\includegraphics[width=0.5\textwidth]{fig/lec01/Automatix_KukaRobot.jpg}
			\caption{Factory robots (source: \href{https://commons.wikimedia.org/wiki/File:Automatix_KukaRobot483.agr.jpg}{Wikimedia Commons}, A.~Reinhold, \href{https://creativecommons.org/licenses/by-sa/4.0/deed.en}{CC BY-SA 4.0})}
		\end{subfigure}
		\hfill
		\begin{subfigure}[b]{0.49\textwidth}
			\centering
			\includegraphics[width=0.5\textwidth]{fig/lec01/electric_drill.jpg}
			\caption{Electric tools (source: \href{https://www.flickr.com/photos/30478819@N08/49940384798}{flickr.com},  M.~Verch, \href{https://creativecommons.org/licenses/by/2.0/}{CC BY 2.0})}
		\end{subfigure}
		\caption*{Examples of electrical machine and drive applications} 
        \label{fig:examples_machine_drives_01}
	\end{figure}
\end{frame}

%%%%%%%%%%%%%%%%%%%%%%%%%%%%%%%%%%%%%%%%%%%%%%%%%%%%%%%%%%%%%
%% Examples of electrical machine and drive applications (2) %%
%%%%%%%%%%%%%%%%%%%%%%%%%%%%%%%%%%%%%%%%%%%%%%%%%%%%%%%%%%%%%
\begin{frame}
	\frametitle{Examples of electrical machine and drive applications (2)}
	% Set up a 2x2 grid of figures
	\begin{figure}
		\ContinuedFloat
		\centering
		\begin{subfigure}[b]{0.49\textwidth}
			\centering
			\includegraphics[width=0.5\textwidth]{fig/lec01/Train.jpg}
			\caption{High-speed trains (source: \href{https://commons.wikimedia.org/wiki/File:Fast_Train_Spain_Class_103_AVE_Siemens_Bridge_Macanet-Massanes.JPG}{Wikimedia Commons}, P.~Elektro, \href{https://creativecommons.org/licenses/by-sa/3.0/deed.en}{CC BY-SA 3.0})}
		\end{subfigure}
		\hfill
		\begin{subfigure}[b]{0.49\textwidth}
			\centering
			\includegraphics[width=0.5\textwidth]{fig/lec01/Electric_Airbus_A320.jpg}
			\caption{Electric aircrafts (source: \href{https://commons.wikimedia.org/wiki/File:Electric_Airbus_A320.jpg}{Wikimedia Commons}, M.~Weinold, \href{https://creativecommons.org/licenses/by-sa/4.0/deed.en}{CC BY-SA 4.0})}
		\end{subfigure}
		\\
		\begin{subfigure}[b]{0.49\textwidth}
			\centering
			\includegraphics[width=0.5\textwidth]{fig/lec01/Pump.jpg}
			\caption{Pumps (source: \href{https://commons.wikimedia.org/wiki/File:Hammelmann_Stationary_unit_with_electric_motor.jpg}{Wikimedia Commons}, Hammelmann, \href{https://creativecommons.org/licenses/by-sa/3.0/deed.en}{CC BY-SA 3.0})}
		\end{subfigure}
		\hfill
		\begin{subfigure}[b]{0.49\textwidth}
			\centering
			\includegraphics[width=0.5\textwidth]{fig/lec01/crane.jpg}
			\caption{Cranes (source: \href{https://commons.wikimedia.org/wiki/File:Hammelmann_Stationary_unit_with_electric_motor.jpg}{Wikimedia Commons}, Belfast Dissenter, \href{https://creativecommons.org/licenses/by-sa/4.0/deed.en}{CC BY-SA 4.0})}
		\end{subfigure}
		\caption*{Examples of electrical machine and drive applications} 
        \label{fig:examples_machine_drives_02}
	\end{figure}
\end{frame}

%%%%%%%%%%%%%%%%%%%%%%%%%%%%%%%%%%%%%%%%%%%%%%%%%%%%%%%%%%%%%
%% Examples of electrical machine and drive applications (2) %%
%%%%%%%%%%%%%%%%%%%%%%%%%%%%%%%%%%%%%%%%%%%%%%%%%%%%%%%%%%%%%
\begin{frame}
	\frametitle{A broad range of nominal power ratings}
	\begin{figure}
		\centering
		\includegraphics[height=0.7\textheight]{fig/lec01/Power_Classes_Examples.pdf}
		\caption{Power range overview (sources: \href{https://www.flickr.com/photos/arthurwolf/5393520058/}{A. Wolf}, \href{https://commons.wikimedia.org/wiki/File:Wald_am_Arlberg-OeBB_Spullersee_power_plant-M1-Rotor-11ASD.jpg}{Asurnipal}, \href{https://www.flickr.com/photos/mouser-nerdbot/7042785635}{M. Williams}, \href{https://de.m.wikipedia.org/wiki/Datei:Stick_blender_Electrolux_AEG_HB_9807_-_stator_of_the_electric_motor-4313.jpg}{R. Spekking}, \href{https://commons.wikimedia.org/wiki/File:Electric_motor_and_transmission_in_a_truck.jpg}{Foxcorner}, \href{https://commons.wikimedia.org/wiki/File:E-bike_electric_motor_shimano_ep_8.jpg}{A.~Tredz} and \href{https://commons.wikimedia.org/wiki/File:2023_Corsair_SP120_RGB_Elite.jpg}{J. Halicki} under varying CC licenses) }
		\label{Power_Classes_Examples}
	\end{figure}
\end{frame}

%%%%%%%%%%%%%%%%%%%%%%%%%%%%%%%%%%%%%%%%%%%%%%%%%%%%%%%%%%%%%
%% Why are electric machines and drives important ? %%
%%%%%%%%%%%%%%%%%%%%%%%%%%%%%%%%%%%%%%%%%%%%%%%%%%%%%%%%%%%%%
\begin{frame}
	\frametitle{Why is knowledge about electric machines and drives important?}
	\begin{varblock}{Electric machines and drives are an essential pillar of the modern society}
		Without electric machines and drives, our todays' society would not be possible. Starting from providing electricity via electrical generators to powering electric vehicles, tools and entire factory production lines, electric machines and drives are everywhere, that is, they enable our today's living standard. 
	\end{varblock}
	\begin{varblock}{Energy efficiency and sustainability is key}
		Electric machines and drives utilize  approx. 50\,\% of the global electricity with about 8 billion electric motors in use in the EU (source: \href{https://commission.europa.eu/energy-climate-change-environment/standards-tools-and-labels/products-labelling-rules-and-requirements/energy-label-and-ecodesign/energy-efficient-products/electric-motors-and-variable-speed-drives_en}{European Commission} and \href{https://iea.blob.core.windows.net/assets/d69b2a76-feb9-4a74-a921-2490a8fefcdf/EE_for_ElectricSystems.pdf}{International Energy Agency}). Therefore, improving their efficiency is an essential factor to reduce the global energy consumption and the associated CO$_2$ emissions.
	\end{varblock}
\end{frame}